Bangla text classifiers meet regional dialects, Bangla typed in Latin letters and posts mixing Bangla, English and
Hindi, and their confidence is usually repaired with temperature scaling, which cannot change a single decision. This
empirical study asks which part of a fine-tuned classifier fails under such shifts and what repairs it. Seven encoders
fine-tuned on Bangla offensive-language and sentiment data, encoders trained on English and Hindi--English sources,
and four instruction-tuned language models are evaluated on 21 source--target pairs from 11 target corpora. The
macro-F1 gain of a bias correction rises with the shift of the predicted-label distribution (Spearman correlation
0.69). In 12 pairs the ranking survives while the decision threshold moves: a label-free bias correction raises
macro-F1 from 0.527 to 0.620, whereas temperature scaling changes no prediction. Fitted jointly with a bias, the
apparent temperature change in these pairs almost vanishes; the class-prior change explains little of the needed
correction, and label-shift estimators fail. The single ranking collapse, on a student-written tri-lingual corpus,
follows the source corpus in a translation factorial but reverses on natural romanised Bangla. A triage procedure
turns these findings into a deployment rule with an explicit label budget.
